\begin{abstract}
% Workflow schedulers must request memory before a task runs. The bytes staged
% as task inputs are available, but they can differ substantially
% from the bytes the program actually consumes through explicit reads and
% memory-mapped page faults. We present Consumption-Aware Memory Prediction
% (\toolname), a feedback loop that selectively audits running tasks with eBPF,
% stores their measured consumed data, and uses only completed observations to
% allocate later tasks. On 33,516 task records from six workflows, we
% test memory allocation first without and then with summaries of data consumed
% by earlier audited tasks and predicted consumption. On 5,026 test tasks, adding consumed-data history
% reduces predicted out-of-memory events by 64.37\%,
% reduces runtime-weighted wastage by 37.18\%, and increases total requested
% memory by 31.50\%. A fixed three-gate policy targets allocation risk,
% consumed-data information, and discovery. At equal audit counts, it finds underallocated
% tasks reliably better than random from 20--50\% budgets. When older
% consumed-data measurements exist, it also finds more tasks worth re-auditing
% than random from 10--50\% budgets. A live Bowtie2 case study completes the
% cold-run, persist, retrain, and fresh-rerun loop. The warm run reduces
% requested memory by 71.96\% while retaining 99.78\% 
% first-attempt coverage. In a separate 21,337-task comparison with Sizey,
% \toolname reduces underallocations by 92.02\% and median absolute percentage
% error by 40.31\%, at the cost of higher memory requests.
% These results show that data consumed by completed tasks is a useful
% allocation signal, while also quantifying its memory cost and the present
% limits of selective auditing.

Distributed scientific workflow tasks exhibit significant variation in memory consumption, making accurate peak memory estimation critical for resource management and avoiding out-of-memory (OOM) failures. Existing workload schedulers rely on static features like file sizes or tool names, which fail to capture non-linear memory scaling. We present Consumption-Aware Memory Prediction (\toolname), a closed-loop framework that utilizes fine-grained, consumption-level provenance—the exact number of bytes a task reads—measured dynamically via eBPF kernel probes. To minimize tracing overhead across parallel task executions at scale, we develop an active-learning-based, multi-gate selective auditing policy that samples only the most informative workflow tasks for model retraining. We evaluate \toolname on thousands of production-scale bioinformatics workflow records executed under Slurm. Compared to size-based scheduling baselines, incorporating dynamic data-consumption features reduces prediction median absolute percentage error by up to 40.31\%. Furthermore, \toolname successfully reduces predicted OOM events by 64.37\% while slashing runtime-weighted memory wastage by 37.18\%. Our selective auditing framework consistently outperforms random sampling across tight profiling budgets, proving that kernel-level consumption data provides a highly scalable and practical feedback signal for high-performance workflow resource allocation.
A live Bowtie2 cold-to-warm evaluation reduces requested memory by 71.96\% while retaining 99.78\% first-attempt coverage.
\end{abstract}

\begin{IEEEkeywords}
scientific workflows, memory prediction, eBPF, selective auditing, Slurm
\end{IEEEkeywords}
